% 07_body.tex — Day2 산출물 2/5 (⑦) 요구사항 추적표 골격 (빈 템플릿 / 완성 예시 공용 본문)
% 드라이버: 07_요구사항추적표_빈템플릿.tex (\wbblanktrue) / 07_요구사항추적표_완성예시.tex (\wbblankfalse)
% 내용 출처: PM트랙/Day2_워크북.md §2.3(항목 가이드) §2.4(빈 템플릿) §2.5(온담 P1 완성 예시본)
\documentclass[10pt,a4paper]{article}
\usepackage[a4paper,margin=16mm,top=14mm,bottom=20mm]{geometry}
\def\DocNum{2}
\def\DocTitle{요구사항 추적표 골격}
\def\DocSub{⑦ Requirements Traceability Matrix (Skeleton)}
\def\DocVariant{\B{빈 템플릿}{온담 P1 완성 예시}}
\usepackage{wbstyle}
\begin{document}

\docheader
 {\Bg{[정식 명칭과 코드]}{차세대 주문·재고 통합 플랫폼 구축 (ONE ONDAM P1)}}
 {\Bg{[RTM v0.x / 베이스라인]}{RTM v0.9 / 2026-03-27 (베이스라인 v1.0, FZ-03)}}
 {\Bg{[PMO 요구사항 담당 실명]}{P1 PMO 요구사항 담당\rd{7mm}}}
 {\Bg{[Senior User 실명 — RACI 5행 A]}{오정근 영업본부장 (RACI 5행 A)\rd{7mm}}}

% ------------------------------------------------------------------ 1
\secbar{1. 문서 정보와 계층 규칙}
\guide{RTM 버전, 기준 요구사항 베이스라인(동결 회차), 계층 정의(HR → StRS → SyRS → SRS)와 층별 건수, ID 규칙, 상위 문서(헌장 §4, 범위 기술서). 건수 검산(층 합 = 전체)을 여기 적어 매 동결마다 갱신한다. 층을 섞으면 — StRS(“한 화면에서 발주할 수 있다”)와 SRS(“수량 입력은 정수만”)가 같은 층에 있으면 — 커버리지 100\%가 무의미하다.}
{\footnotesize
\begin{tabularx}{\linewidth}{|L{30mm}|Y|L{22mm}|Y|}\hline
\hd{항목} & \hd{내용} & \hd{항목} & \hd{내용} \\ \hline
\B{\rd{9mm}}{}RTM 버전 / 베이스라인 & \Bg{[v0.x / 요구사항 베이스라인 vN, 동결 회차 FZ-nn, 일자]}{v0.9 / 요구사항 베이스라인 v1.0 (FZ-03, 2026-03-25 동결, S6 종료)} & 계층·건수 & \Bg{[HR\ \ / StRS\ \ / SyRS\ \ / SRS\ \ = 합계]}{HR 8 / StRS 214 / SyRS 396 / SRS 570 = 1,180 (P-09)} \\ \hline
\B{\rd{9mm}}{}ID 규칙 & \Bg{[HR-nn, StRS-nnn, SyRS-nnn, SRS-nnn, TC-]}{HR-nn, StRS-nnn(HR별 20번 단위), SyRS-nnn, SRS-nnn, TC-<영역>-nnn} & 상위 문서 & \Bg{[헌장 §4, 범위 기술서 vN]}{헌장 v1.0 §4·§13, 범위 기술서 v0.9} \\ \hline
\B{\rd{9mm}}{}작성 / 승인 & \Bg{[PMO 담당 / Senior User 실명]}{P1 PMO(요구사항 담당) / 오정근 Senior User (RACI 5행 A)} & 감리 기준 & \Bg{[게이트별 연결률 목표]}{G2 설계·시험 연결 100\%, G3 검수 연결 100\% (헌장 §3)} \\ \hline
\end{tabularx}}

\vspace{2mm}\noindent\textbf{\small 분류 규칙 (MoSCoW)} \guide{Must = 규제·가용성·헌장 고정 항목·편익 직접 기여 / Should = 편익 간접 기여, 이연 시 편익 지연 / Could = 편의 기능 / Won't(이번 릴리스) = 헌장 제외 범위 또는 타 프로젝트 소관. \textbf{Must 상한은 릴리스별 FP 기준}으로 정한다 — 전부 Must면 우선순위가 없는 것이고 헌장 §13의 조정 변수(범위)가 작동하지 않는다. 우선순위는 Senior User가 정하고 PM은 규칙을 지킨다.}
\begin{fieldbox}
\ifwbblank
\g{Must = \hrulefill}\par\vspace{3mm}\g{Should = \hrulefill}\par\vspace{3mm}\g{Could = \hrulefill}\par\vspace{3mm}\g{Won't = \hrulefill}\par\vspace{3mm}\g{Must 상한 \hspace{8mm}\% (FP 기준) — 초과 시 조정자: \hrulefill}\par
\else
Must = 규제·가용성·헌장 §13 고정 항목·편익 직접 기여(B-04·B-05·B-07·B-08) / Should = 편익 간접 기여 또는 이연 시 편익 지연 / Could = 편의·리포트 / Won't(이번 릴리스) = 헌장 제외 범위 또는 P2·P3·P5 소관. \textbf{Must 상한 = 릴리스별 FP의 60\%.} 초과 시 오정근이 Should로 내릴 항목 지정.
\fi
\end{fieldbox}

% ------------------------------------------------------------------ 2 (가로)
\landscapeon
\secbar{2. 상위 요구 노드 (HR)}
\guide{헌장 §4 상위 요구를 ID·문장(헌장 그대로)·요구 제공자·유형(기능/비기능/규제)·연결 편익(B-ID)으로. StRS는 반드시 HR 하나에 매달린다. HR을 새로 쓰거나(헌장과 다른 문장), HR에 매달리지 않는 StRS를 “기타”로 받지 말 것 — 기타는 범위 밖이거나 헌장 수정 요청이다. 연결 편익 열이 우선순위 판정의 근거다.}
{\footnotesize
\begin{longtable}{|L{14mm}|L{130mm}|L{38mm}|L{24mm}|L{48mm}|}\hline
\hd{HR} & \hd{요구 문장 (헌장 §4 그대로)} & \hd{요구 제공자 (S-ID 실명)} & \hd{유형} & \hd{연결 편익·제약} \\ \hline \endhead
\ifwbblank
\rd{11mm}HR-01 & \g{[헌장 §4 문장]} & \g{[S-nn 실명]} & \g{[기능/비기능/규제]} & \g{[B-nn / 제약 n]} \\ \hline
\rd{11mm}HR-02 & & & & \\ \hline
\rd{11mm}HR-03 & & & & \\ \hline
\rd{11mm}HR-04 & & & & \\ \hline
\rd{11mm}HR-05 & & & & \\ \hline
\rd{11mm}HR-06 & & & & \\ \hline
\rd{11mm}HR-07 & & & & \\ \hline
\rd{11mm}HR-08 & & & & \\ \hline
\else
HR-01 & 매장·물류센터·공장 출하 재고를 단일 원장에서 준실시간(15분 이내)으로 조회할 수 있다 & 박세연(S-04), 한동석(S-05) & 기능 & B-02 재고 정확도, B-04 폐기율 \\ \hline
HR-02 & 자사앱·배달3사·POS·가맹 발주·B2B EDI의 주문이 하나의 주문 원장에 기록되고 매장 간 재고 이동·픽업이 가능하다 & 오정근(S-03) & 기능 & B-01 온라인 비중, B-05 배달수수료 회피 \\ \hline
HR-03 & 가맹점주가 포털에서 발주·마감·출고 확인·배송 추적을 한 화면에서 처리할 수 있다 & 서정필(S-10), 오정근(S-03) & 기능 & B-08 리드타임, B-03 결품률 \\ \hline
HR-04 & 성수기(7~8월·12월·명절 전 2주)에 매장 시스템 변경이 발생하지 않도록 배포 창이 설계되어 있다 & 오정근(S-03) & 비기능·운영 & 헌장 제약 1 \\ \hline
HR-05 & 운영 가용성 99.5\% 이상, 월 허용 다운타임 3시간 39분 이내, 컷오버 시 롤백 절차 보유 & 박준영(S-08) & 비기능 & O-18 SLA, 헌장 §14 \\ \hline
HR-06 & 개인정보 처리항목 142개·위탁사 19개·재위탁 4건이 설계에 반영되고 망분리·전자금융 요건을 충족한다 & 최영주(S-07) & 비기능·규제 & 규제 기한 2026-09-11, 2027-03-31 \\ \hline
HR-07 & IT본부 47명이 벤더 없이 운영·변경할 수 있는 아키텍처와 문서(인터페이스 47본 설계문서 포함)를 갖춘다 & 박세연(S-04) & 비기능 & 이사회 승인 조건 5, 2023 감사 지적 \\ \hline
HR-08 & SAP ECC 6.0과의 연동은 2027-12 이후 마이그레이션을 고려한 구조(애드온·API 분리)로 설계한다 & 박세연(S-04) & 비기능 & 헌장 제외 1 (SAP 교체) \\ \hline
\fi
\end{longtable}}

% ------------------------------------------------------------------ 3~9 (가로)
\clearpage
\secbar{3~9. 추적 행 — StRS ID · 요구 문장 · 제공자 · 우선순위 · 하위 추적 · 설계 · 시험 · 검수 · 상태}
\guide{한 행 = StRS 한 건. 문장은 “~할 수 있다 / ~이어야 한다”의 판정 가능한 문장(해결책 “Kafka로 구현한다”는 요구가 아니다). 제공자는 부서가 아니라 사람(S-ID) — “현업”이라고 적힌 행이 동결 직전에 들어온 요구의 대부분이다. 하위 추적은 “SyRS n / SRS n (SRS-nnn~nnn)”, 건수 0이면 아직 설계되지 않은 요구. 시험 열에 비기능 요구도 시험 ID를 적는다(“검토”는 시험이 아니다). 검수 열은 “인수 기준 n / Gn / 실명” — \textbf{비어 있으면 G3에서 커버리지 100\%를 선언할 수 없고, 인수자가 PM이면 그 요구는 아직 주인을 찾지 못한 것이다.} 기각도 이력이다 — 지우지 않는다.}
{\scriptsize\setlength{\tabcolsep}{2pt}
\begin{longtable}{|L{11mm}|L{42mm}|L{9mm}|L{15mm}|L{16mm}|L{11mm}|L{12mm}|L{15mm}|L{18mm}|L{12mm}|L{19mm}|L{22mm}|L{11mm}|L{22mm}|}\hline
\hd{StRS} & \hd{요구 문장} & \hd{HR} & \hd{제공자 (S-ID)} & \hd{접수 경로·일자} & \hd{MoSCoW} & \hd{릴리스} & \hd{SyRS / SRS (대표 ID)} & \hd{설계 산출물} & \hd{WBS} & \hd{시험 케이스} & \hd{검수 / 게이트 / 인수자} & \hd{상태 (FZ)} & \hd{변경 이력} \\ \hline \endhead
\ifwbblank
\rd{11mm}\g{StRS-[\ ]} & \g{[판정 가능한 한 문장]} & \g{HR-[\ ]} & \g{[S-nn 실명]} & \g{[경로, 일자]} & \g{[M/S/C/W]} & \g{[R1~R5]} & \g{[n / n (ID~ID)]} & \g{[문서 §]} & \g{[1.n.m]} & \g{[TC-]} & \g{[기준 n / Gn / 실명]} & \g{[상태 (FZ-nn)]} & \g{[CR-nn 일자 내용 ΔFP 승인자]} \\ \hline
\rd{11mm} & & & & & & & & & & & & & \\ \hline
\rd{11mm} & & & & & & & & & & & & & \\ \hline
\rd{11mm} & & & & & & & & & & & & & \\ \hline
\rd{11mm} & & & & & & & & & & & & & \\ \hline
\rd{11mm} & & & & & & & & & & & & & \\ \hline
\rd{11mm} & & & & & & & & & & & & & \\ \hline
\rd{11mm} & & & & & & & & & & & & & \\ \hline
\rd{11mm} & & & & & & & & & & & & & \\ \hline
\rd{11mm} & & & & & & & & & & & & & \\ \hline
\rd{11mm} & & & & & & & & & & & & & \\ \hline
\rd{11mm} & & & & & & & & & & & & & \\ \hline
\rd{11mm} & & & & & & & & & & & & & \\ \hline
\rd{11mm} & & & & & & & & & & & & & \\ \hline
\rd{11mm} & & & & & & & & & & & & & \\ \hline
\rd{11mm} & & & & & & & & & & & & & \\ \hline
\rd{11mm} & & & & & & & & & & & & & \\ \hline
\rd{11mm} & & & & & & & & & & & & & \\ \hline
\rd{11mm} & & & & & & & & & & & & & \\ \hline
\rd{11mm} & & & & & & & & & & & & & \\ \hline
\else
StRS-031 & 매장 판매 즉시 중앙 재고 원장이 차감되며 반영 지연은 15분 이내다 & HR-01 & S-04 박세연 & 설계 검토회 S2, 2026-01-28 & Must & R2 & 4 / 7 (SRS-201~207) & 아키텍처 정의서 §4.2 & 1.3.1, 1.3.2 & 통합 TC-INV-001~010, 성능 TC-PERF-03 & 인수 기준 ③-① / G3 / 한동석 & 승인 (FZ-01) & — \\ \hline
StRS-032 & 이천센터 입출고가 대성 WMS로부터 준실시간으로 재고 원장에 반영되고, 연동 불가 시 야간 배치 2회로 자동 전환된다 & HR-01 & S-05 한동석 & 물류본부 협의 S3, 2026-02-11 & Must & R2–R3 & 3 / 6 (SRS-208~213) & 인터페이스 설계서 IF-3PL-01~05 & 1.3.3 & TC-INV-021~040 & 인수 기준 ③-② / G3 / 한동석 & 승인 (FZ-02) & — \\ \hline
StRS-033 & 안성·이천 공장 출하 재고가 온담푸드시스템 계약 범위 안에서 재고 원장에 연동된다 & HR-01 & S-05 한동석, S-14 조성태 & 서면 검토 S4, 2026-02-25 & Should & R3 & 2 / 4 (SRS-214~217) & IF-FAC-01~02 & 1.3.5 & TC-INV-041~046 & 인수 기준 ③-① / G3 / 조성태·박세연 & 승인 (FZ-02) & — \\ \hline
StRS-034 & 재고 정확도 리포트는 순환실사 SKU(69\%)와 전체 SKU(100\%) 두 분모를 병기한다 & HR-01 & S-06 나영선, S-05 한동석 & 분모 협의 S5, 2026-03-10 & Must & R3 & 2 / 3 (SRS-218~220) & 리포트 정의서 RPT-INV-02 & 1.3.2 & TC-INV-050~052 & 인수 기준 ③-③ / G3 / 한동석 & 승인 (FZ-03) & — \\ \hline
StRS-051 & 5개 채널의 주문이 단일 주문 원장에 채널 식별자와 함께 기록된다 & HR-02 & S-03 오정근 & 설계 검토회 S2, 2026-01-28 & Must & R3 & 5 / 9 (SRS-301~309) & 주문 원장 모델 v1.0 §2 & 1.2.1, 1.2.2 & TC-ORD-001~030 & 인수 기준 ②-① / G3 / 오정근 & 승인 (FZ-01) & — \\ \hline
StRS-052 & 앱 주문을 지정 매장에서 픽업할 수 있고, 매장은 픽업 주문을 POS에서 처리한다 & HR-02 & S-03 오정근 & 설계 검토회 S3, 2026-02-11 & Must & R4 & 3 / 6 (SRS-310~315) & 옴니 시나리오 정의서 OMN-01 & 1.2.5 & TC-ORD-060~075, UAT-ORD-03 & 인수 기준 ②-② / G3 / 오정근 & 승인 (FZ-02) & — \\ \hline
StRS-053 & 매장 재고 부족 시 인근 매장 재고를 시스템이 자동 제안하고 이동을 예약한다 & HR-02 & S-03 오정근 & 설계 검토회 S4, 2026-02-25 & Should & \textbf{R5} (이연 가능, 헌장 §13) & 2 / 5 (SRS-316~320) & OMN-02 & 1.2.5 & TC-ORD-076~085 & 인수 기준 ②-② / G4 / 오정근 & 승인 (FZ-02) & — \\ \hline
StRS-054 & 배달3사 주문을 가맹점에 자동 배분하고 가맹 매출 몫을 정산한다 & HR-02 & S-10 서정필 & 협의회 인터뷰 S4, 2026-02-26 & \textbf{Won't} & — & 0 / 0 & — & — & — & — & 기각 (FZ-03) & CR-05 2026-03-25 Won't 판정 — 가맹 매출 몫은 부속서 협상(P3) 소관, 오정근 승인. 기록 유지 \\ \hline
StRS-071 & 가맹점주가 발주·마감·출고 확인·배송 추적을 한 화면(웹·모바일)에서 처리한다 & HR-03 & S-10 서정필 & 협의회 인터뷰 S1, 2026-01-14 & Must & R3 (선행 릴리스 2026-Q4) & 4 / 8 (SRS-401~408) & 포털 화면 정의서 FRN-01~06 & 1.4.1, 1.4.2 & UAT-FRN-001~030 & 인수 기준 ④-① / 선행 릴리스·G3 / 오정근·서정필 & 승인 (FZ-01) & — \\ \hline
StRS-072 & 기존 엑셀 발주 양식을 업로드하면 발주로 변환된다 & HR-03 & S-03 오정근 & 매장운영 검토 S3, 2026-02-12 & Should & R3 & 2 / 4 (SRS-409~412) & FRN-07 & 1.4.2 & UAT-FRN-031~038 & 인수 기준 ④-① / G3 / 오정근 & 승인 (FZ-03) & CR-03 2026-03-11 양식 3종→5종 확장, ΔFP +22, 허용범위 내 P1 PM 승인, CCB 사후 보고 \\ \hline
StRS-073 & 발주 접수부터 출고 확정까지 12시간 이내가 되도록 마감·출고 확정 흐름이 설계된다 & HR-03 & S-03 오정근 & 설계 검토회 S2, 2026-01-29 & Must & R3 & 3 / 5 (SRS-413~417) & 발주 흐름 정의서 FRN-FLOW-01 & 1.4.2, 1.5.4 & TC-FRN-040~048, UAT-FRN-039 & 인수 기준 ④-② / G3 / 오정근 & 승인 (FZ-01) & — \\ \hline
StRS-091 & 매장 SW 배포는 성수기 창(7~8월·12월·명절 전 2주) 밖에서만 실행되도록 배포 캘린더에 잠금이 있다 & HR-04 & S-03 오정근 & 설계 검토회 S1, 2026-01-15 & Must & R2 & 2 / 3 (SRS-501~503) & 배포 절차서 DEP-01 & 1.1.8, 1.6.9 & TC-DEP-001~006 & 인수 기준 ①-③ / G3 / 오정근 & 승인 (FZ-01) & — \\ \hline
StRS-092 & 가맹 단말 SW 배포는 매장별로 선택·예약 실행할 수 있고 실패 시 이전 버전으로 자동 복귀한다 & HR-04 & S-03 오정근, S-10 서정필 & 협의회 인터뷰 S1, 2026-01-14 & Must & R4 & 3 / 5 (SRS-504~508) & DEP-02 & 1.1.6, 1.1.8 & TC-POS-150~180 & 인수 기준 ①-③·④ / 컷오버 / 오정근 & 승인 (FZ-01) & — \\ \hline
StRS-101 & 플랫폼 가용성은 월 99.5\% 이상이며 이중화·장애 전환이 자동이다 & HR-05 & S-08 박준영 & 운영 요건 검토 S2, 2026-01-30 & Must & R1 (기반) & 4 / 6 (SRS-601~606) & 아키텍처 정의서 §6 가용성 & 1.5.1, 1.5.6 & 장애 전환 TC-DR-001~012 & 인수 기준 ⑤-① / G3, 컷오버 후 3개월 / 박준영 & 승인 (FZ-01) & — \\ \hline
StRS-102 & 컷오버 실패 시 6영업일 이내 OD-POS로 롤백할 수 있는 절차와 도구가 있다 & HR-05 & S-08 박준영 & 운영 요건 검토 S2, 2026-01-30 & Must & R4 & 2 / 4 (SRS-607~610) & 롤백 절차서 RB-01 & 1.1.8, 1.6.9 & 롤백 리허설 기록 RH-01~02 & 인수 기준 ①-④ / G3 / 박준영·정미란 & 승인 (FZ-01) & — \\ \hline
StRS-103 & 장애·성능 지표가 APM에서 실시간 조회되고 서비스데스크 11명이 운영 문서만으로 1차 대응한다 & HR-05 & S-08 박준영 & 운영 요건 검토 S4, 2026-02-27 & Should & R5 & 3 / 5 (SRS-611~615) & 운영 기술 문서 OPS-01~04 & 1.5.6, 1.6.10 & 운영 리허설 TC-OPS-001~010 & 인수 기준 ⑧-② / G4 / 박준영 & 승인 (FZ-03) & — \\ \hline
StRS-121 & 개인정보 처리항목 142개가 수집·보관·파기 주기와 함께 설계에 반영된다 & HR-06 & S-07 최영주 & 규제 검토 의견 S3, 2026-02-13 & Must & R2 (2026-09-11 전) & 6 / 9 (SRS-701~709) & 규제 요건 매핑표 REG-01 & 1.5.1, 1.6.6 & TC-REG-001~142 & 인수 기준 ⑤-① / 규제 ① 2026-09-11 / 최영주 & 승인 (FZ-02) & — \\ \hline
StRS-122 & 위탁사 19개·재위탁 4건의 처리 위탁 구조가 시스템 접근 권한에 반영되며, 별도 법인 처리 주체가 구분된다 & HR-06 & S-07 최영주 & 규제 검토 의견 S3, 2026-02-13 & Must & R2 & 3 / 5 (SRS-710~714) & REG-02, 권한 모델 AUTH-01 & 1.5.1, 1.3.5 & TC-REG-143~160 & 인수 기준 ⑤-① / 규제 ① / 최영주 & 승인 (FZ-02) & — \\ \hline
StRS-123 & 앱 선불충전 관련 거래는 전자금융 요건(망분리·거래 기록 보존)을 충족한다 & HR-06 & S-07 최영주 & 규제 검토 의견 S5, 2026-03-12 & Must & R3 (신청서 2027-01-29 전) & 3 / 6 (SRS-715~720) & REG-03, 망분리 구성도 & 1.5.6, 1.2.3 & 보안진단 2회차, TC-REG-161~175 & 인수 기준 ⑤-① / 규제 ② 2027-03-31 / 최영주 & 승인 (FZ-03) & — \\ \hline
StRS-141 & 인터페이스 47본의 설계문서가 표준 양식으로 존재하고 문서 부재 13본이 복원된다 & HR-07 & S-04 박세연 & 설계 검토회 S1, 2026-01-15 & Must & R0~R1 (복원 S2–S6) & 2 / 3 (SRS-801~803) & IF-001~047 & 1.5.5 & 감리 문서 검토 G2 & 인수 기준 ⑥ / G2 / 박세연 & 승인 (FZ-01) & — \\ \hline
StRS-142 & 소스·빌드·배포 절차가 발주사 저장소에서 클린 빌드로 재현된다 & HR-07 & S-04 박세연 & 형상관리 검토 S2, 2026-01-29 & Must & R5 & 2 / 4 (SRS-804~807) & 형상관리 계획 CM-01 & 1.6.2, 1.6.10 & 재현 시험 TC-CM-001~005 & 인수 기준 ⑦ / G4 / 박세연 & 승인 (FZ-01) & — \\ \hline
StRS-143 & 상용SW(API GW·MDM·APM)의 설정·정책 변경을 IT본부가 벤더 없이 수행할 수 있는 관리 화면·문서가 있다 & HR-07 & S-04 박세연 & 설계 검토회 S4, 2026-02-26 & Should & R5 & 3 / 5 (SRS-808~812) & OPS-05~07 & 1.5.2, 1.5.3, 1.6.10 & TC-OPS-011~020 & 인수 기준 ⑧-① / G4 / 박세연·박준영 & 승인 (FZ-02) & — \\ \hline
StRS-161 & SAP ECC 연동은 애드온과 API를 분리해 ECC 교체 시 API 계층만 유지되도록 설계된다 & HR-08 & S-04 박세연 & 설계 검토회 S2, 2026-01-29 & Must & R2 & 4 / 6 (SRS-901~906) & 아키텍처 정의서 §5 SAP 연동 & 1.5.4 & TC-INT-001~040 & 인수 기준 ⑤-① / G3 / 박세연 & 승인 (FZ-01) & — \\ \hline
StRS-162 & OD-POS↔SAP 11본·가맹발주↔SAP 4본이 신 플랫폼 API로 대체되고 병행 기간 중 양방향 정합이 유지된다 & HR-08 & S-04 박세연 & 설계 검토회 S3, 2026-02-12 & Must & R3 & 3 / 5 (SRS-907~911) & IF-SAP-01~15 & 1.5.4, 1.5.5 & TC-INT-041~070 & 인수 기준 ⑤-① / G3 / 박세연 & 승인 (FZ-02) & — \\ \hline
StRS-163 & SAP 마스터 변경이 MDM에 15분 이내 반영되고 역방향 반영은 승인 절차를 거친다 & HR-08 & S-04 박세연, S-15 강현도 & 구매팀 검토 S5, 2026-03-11 & Could & R3 & 1 / 2 (SRS-912~913) & MDM 동기화 정의서 MDM-03 & 1.5.3, 1.5.4 & TC-INT-071~075 & 인수 기준 ⑤-① / G3 / 박세연 & 승인 (FZ-03) & — \\ \hline
\fi
\end{longtable}}
\B{\small\g{(계속되면 별지 — 같은 열 순서로)}\par}{\small 샘플 25건 = HR-01 4 + HR-02 4(StRS-054 Won't 포함) + HR-03 3 + HR-04 2 + HR-05 3 + HR-06 3 + HR-07 3 + HR-08 3. 전체 214건 중 구조가 보이는 골격 — 나머지는 같은 열로 늘린다.\par}

% ------------------------------------------------------------------ 10 (가로)
\secbar[70mm]{10. 커버리지 집계}
\guide{층별 건수, 상위 연결률(StRS 중 HR에 매달린 비율), 하위 분해율(SyRS·SRS 1건 이상), 시험 연결률, 검수 연결률, MoSCoW 비율. 동결마다 갱신. 감리 G2 “설계·시험 연결 100\%”, G3 “검수 연결 100\%”가 목표. 집계를 감리 직전에 처음 하지 말 것. \textbf{Must 비율은 건수가 아니라 FP 가중}으로 본다.}
{\footnotesize
\begin{tabularx}{\linewidth}{|L{40mm}|C{16mm}|C{22mm}|C{34mm}|C{30mm}|C{30mm}|Y|C{34mm}|}\hline
\hd{층} & \hd{건수} & \hd{상위 연결률} & \hd{하위 분해율} & \hd{시험 연결률} & \hd{검수 연결률} & \hd{MoSCoW (건수)} & \hd{Must \% (FP 가중)} \\ \hline
\ifwbblank
\rd{9mm}\g{샘플 StRS} & & \g{[\ ]\%} & \g{[\ ]\%} & \g{[\ ]\%} & \g{[\ ]\%} & \g{M\ \ / S\ \ / C\ \ / W\ \ } & \g{[\ ]\%} \\ \hline
\rd{9mm}\g{전체 StRS} & & & & & & & \\ \hline
\rd{9mm}SyRS & & & — & & — & — & — \\ \hline
\rd{9mm}SRS & & & — & & — & — & — \\ \hline
\else
샘플 StRS & 25 & 100\% & 96\% (Won't 1건 제외 시 100\%) & 100\% (Won't 제외) & 100\% (Won't 제외) & M 14 / S 7 / C 3 / W 1 & 약 58\% [추가 설정] \\ \hline
전체 StRS (S6 종료) & 214 & 100\% & 71\% & 43\% & 0\% (G3까지 채움) & M 131 / S 58 / C 19 / W 6 & 61\% → G1까지 60\% 이하로 조정(오정근) \\ \hline
SyRS & 396 & 100\% & — & 38\% & — & — & — \\ \hline
SRS & 570 & 100\% & — & 35\% & — & — & — \\ \hline
\fi
\end{tabularx}}
\B{}{\small 검산: 214 + 396 + 570 = 1,180 (P-09). 범위 허용범위 ±3\% = 1,180 × 0.03 = 35.4 → ±35건(헌장 §12). FP 기준 ±3\% = 12,400 × 0.03 = 372점.\par}
\landscapeoff
\end{document}
